% Placeholder 5-paragraph cover letter demonstrating the mandatory structure and >=60% page-height rule.
\vspace{0.4em}
\today

\vspace{0.4em}
\noindent
Hiring Team\\
Example Company\\
Example City, Country

\vspace{0.3em}
\LetterSubject{Application for Example Data \& AI Engineering Internship (REQ-0000)}

Dear Hiring Team,

% Paragraph 1: Hook Opening (why this company, team, and problem space — no generic opener)
Building reliable operational intelligence where automated data pipelines meet day-to-day engineering decisions is the exact problem space I want to focus on during my MSc internship, which is why the Example Data \& AI Engineering Internship at Example Company immediately stood out to me. Seeing how your team bridges raw telemetry and internal engineering workflows into production decision tools aligns directly with my focus on shipping maintainable, end-to-end data and software systems rather than isolated notebook prototypes.

% Paragraph 2: Focused Resume Proof (1-2 concrete anchors, not a full resume recap)
In my current internship at Example Enterprise Tech Corp, I evolved a set of fragmented Python and SQL reporting scripts into a containerized business operational web application built with FastAPI, PostgreSQL, and Docker for roughly thirty regional stakeholders, taking the system through formal architecture and security reviews. Earlier at Example Global HealthTech Co., I tackled the upstream side of operational reliability by investigating master-to-transactional data discrepancies across supply chain workflows in SQL and Python and turning those findings into automated reconciliation checks and Power BI monitoring views.

% Paragraph 3: Company & Role Understanding (specific observation on workflow/challenges from JD)
What makes this role at Example Company particularly compelling is the emphasis in your job description on connecting automated data quality validation with actionable developer and operations tooling. In many organizations, analytics dashboards and operational execution live in separate silos: metrics surface anomalies days later, but teams still rely on manual spreadsheet triage to resolve them. Closing that loop requires engineers who care as much about clean data contracts, schema evolution, and reproducible pipelines as they do about the usability of the interface placed in front of domain teams.

% Paragraph 4: Forward-Looking First-Weeks Execution Plan (Mandatory depth & onboarding plan)
If I join the team, my priority in the first two to three weeks would be to map the existing ingestion and validation pipeline end to end—auditing how raw operational events are transformed, where silent schema drift or manual overrides currently occur, and which recurring bottlenecks cost engineers the most triage time. From there, I would focus on delivering a well-tested, modular slice of the pipeline—whether tightening automated data contract checks in Python and SQL or exposing a lightweight operational endpoint—while using structured AI-assisted engineering workflows to iterate rapidly on documentation and regression tests alongside senior engineers.

% Paragraph 5: Logistics & Specific Closing
I am currently enrolled in the MSc program at Example University of Technology and am available to start in [Month Year] for five to six months on a full-time basis ([4--5 days per week]) in Example City. I would welcome the opportunity to discuss how my experience building production-bound Python and PostgreSQL operational tools can support your team's upcoming roadmap.

\vspace{0.6em}
Sincerely,\\[0.8em]
[Full Name]
